\chapter{Background}

\section{Related Work}

Several studies have investigated applying Dynamic Time Warping (DTW)\todo{remove when explained earlier} on spatial domains and 
weather data separately \todo{add citations}, and other studies have investigated the decomposition
of weather forecast error into spatial and temporal components \todo{add citations}. The intersection between these 
topics, however, is largely unexplored.

Dynamic Time Warping applied to spatial orderings has been used with proven results in recent
litterature on automatic correlation between well-logs for geological analysis\cite{well-log-correlation}. 
Nilo et al. use the DTW algorithm to correlate features in wells where measurements
are ordered by depth, yielding promising results as an aid in research though still requiring some supervision
in practical deployments \cite{well-log-correlation}. Le Coz et al. use DTW together with spatial 
image warping to correct estimates of total rainfall during events such as convective rainstorms in sub-Saharan Africa.
By warping satellite-based estimates to fit measurements from local weather station networks, they note significant 
improvements over pure satellite-based estimates \cite{gauge-adjustment}. Other notable applications include using DTW
to compare cloud-top temperature and other satellite-based estimation technique time-series in precipitation estimation. \cite{cloudtops,andes-rainfall} 

Decomposing forecast errors into spatial and temporal components is generally first accredited to Hoffman et al. in their 
1995 paper "Distortion Representation of Forecast Errors" \cite{hoffman1995,spatiotemporal-forecast-partition}




















\section{Clustering}

Clustering is the process of partitioning a set of data objects into groups referred to as clusters, that are similar to each other or dissimilar to objects in other clusters~\cite{datamining}.
This is an unsupervised machine learning technique used to explore the dataset, for preprocessing or as a subroutine for further study~\cite{kshape}. 
The two main advantages of clustering is that since it is unsupervised it can easily be applied to real-world datasets that don't need to be curated or processed first, and that it might let us discover groupings that humans would not intuitively think of.
